% ==========================================================
% 奇偶交替条件下的相邻概率
% 难度：★★★ 难题
% ==========================================================

\newcommand{\qProbPermAltParityAdjProbStem}{%
用 $1,2,3,4,5,6$ 组成无重复数字的六位数。在任意相邻两个数字奇偶性不同的条件下，求 $1$ 和 $2$ 相邻的概率。
}

\newcommand{\qProbPermAltParityAdjProbAnswer}{%
$\frac{5}{9}$
}

\newcommand{\qProbPermAltParityAdjProbSolution}{%
\textbf{第一步：确定条件样本空间}\\
奇数 $1,3,5$，偶数 $2,4,6$。\\
奇偶交替的位置类型只能是：\\
“奇偶奇偶奇偶”或“偶奇偶奇偶奇”。\\
每种位置中：奇数有 $3!$ 种排列，偶数有 $3!$ 种。\\
条件样本空间大小为 $\lvert\Omega\rvert = 2\cdot3!\cdot3! = 72$。

\textbf{第二步：计算 $1$ 和 $2$ 相邻的情况}\\
固定一种奇偶位置类型。\\
六个位置中有五对相邻位置：\\
$(1,2),(2,3),(3,4),(4,5),(5,6)$。\\
选择一对放置 $1,2$，有 $5$ 种选择。\\
由于 $1$ 是奇数、$2$ 是偶数，\\
它们的方向由该位置的奇偶性自动确定，\\
不再需要乘 $2$。\\
剩余奇数排有 $2!$ 种，剩余偶数排有 $2!$ 种。\\
固定一种位置类型时：$5\cdot2!\cdot2!=20$ 种。\\
两种奇偶位置类型共有 $40$ 种。

所求概率为 $P = \frac{40}{72} = \frac{5}{9}$。
}

\newcommand{\qProbPermAltParityAdjProbTeachingNotes}{%
\textbf{【避坑与边界说明】}
\begin{itemize}
    \item 分母不是 $6!$，而是满足条件的 $72$ 个数。\\
    概率计算必须统一在条件样本空间内。
    \item $1$ 和 $2$ 在选定相邻位置中的方向\\
    由奇偶位置自动决定，不能再乘 $2$。
    \item 强行“捆绑 $1,2$”也需服从原奇偶结构。\\
    把捆绑当普通对象任意排会使思路变复杂。
\end{itemize}
}
